% =====================================================================
% Sekcja 3 - Metodologia (wersja kompaktowa, ~1-1.5 strony)
% =====================================================================

\section{Metodologia}
\label{sec:metodologia}

\subsection{Pytanie badawcze}
\label{sec:rq}

W tej pracy sprawdzane jest, czy perplexity i pokrewne sygnały
liczone z rozkładu prawdopodobieństw tokenów - średnia entropia
tokenowa, top-1 margin, top-1 mass - przewidują poprawność
odpowiedzi LLM na konkretne pytanie. Pytanie to stawiane jest
w dwóch ustawieniach: w czystym chain-of-thought (CoT) oraz
w agencie ReAct z pojedynczym narzędziem.

Drugie, pomocnicze pytanie dotyczy tego, gdzie wewnątrz trajektorii
agenta najlepiej mierzyć tę pewność: czy mocniejszy sygnał płynie
z segmentów \emph{action} (wywołań narzędzia), \emph{thought}
(rozumowania) czy \emph{answer} (finalnej odpowiedzi).

\subsection{Zbiór danych}
\label{sec:gsm8k}

Eksperymenty przeprowadzono na zbiorze GSM8K~\cite{cobbe2021gsm8k} -
8500 zadań tekstowych z matematyki szkolnej (7473 trenujących
i 1319 testowych), z których każde wymaga od dwóch do ośmiu kroków
rozumowania. Każde zadanie kończy się jednoznaczną odpowiedzią
liczbową poprzedzoną sekwencją \texttt{\#\#\#\#}, dzięki czemu
ocena poprawności sprowadza się do porównania dwóch liczb - bez
potrzeby angażowania drugiego modelu jako sędziego. Przykład
z części testowej:

\begin{quote}
\textit{Natalia sprzedała spinacze 48 swoim przyjaciółkom w kwietniu,
a w maju sprzedała ich o połowę mniej. Ile spinaczy łącznie
sprzedała w kwietniu i maju?}\\
\textit{Rozwiązanie}: $48/2 = 24$ spinaczy w maju; $48 + 24 = 72$
spinaczy łącznie. \texttt{\#\#\#\# 72}
\end{quote}

GSM8K jest dla tego projektu wygodny z trzech powodów. Po pierwsze,
ocena jest zerojedynkowa i obiektywna - liczba albo się zgadza,
albo nie. Po drugie, zadania są krótkie i mieszczą się w oknie
kontekstowym małych modeli (0,5-1,5B parametrów). Po trzecie,
modele tej skali osiągają na GSM8K dokładność rzędu 20-60\%,
czyli ani nie zgadują na ślepo, ani nie rozwiązują wszystkiego -
co jest pożądane przy badaniu korelacji pewności z poprawnością.
GSM8K jest też standardowym benchmarkiem dla prac o rozumowaniu
krok po kroku od~\cite{wei2022cot}.

Z części testowej losowano 50 zadań (faza pilotażowa) lub 60 zadań
(eksperymenty główne) przy ustalonych ziarnach. Modele nie były
dotrenowywane - wszystkie wyniki pochodzą z czystej ewaluacji.

\subsection{Modele}
\label{sec:modele}

W eksperymentach użyto dwóch modeli z tej samej rodziny, różniących
się tylko skalą: Qwen2.5-0.5B-Instruct (w fazie pilotażowej) oraz
Qwen2.5-1.5B-Instruct (w eksperymentach głównych). Pozwala to
sprawdzić, czy obserwacje z fazy pilotażowej utrzymują się przy
przejściu na większy model. Spełnia to też wymóg co najmniej dwóch
modeli w ramach kursu Warsztaty AI.

\subsection{Dekodowanie i skoring}
\label{sec:dekodowanie}

We wszystkich eksperymentach stosowane jest dekodowanie zachłanne
(\emph{greedy decoding}) - na każdym kroku wybierany jest token
o najwyższym prawdopodobieństwie. Eliminuje to szum samplingu
i sprawia, że wartości PPL między pytaniami można porównywać.
Ma to swoją cenę: ponieważ model zawsze wybiera swój top-1 token,
jego prawdopodobieństwa są systematycznie wysokie, a wartości PPL
skupiają się blisko 1,0. Jest to jedno z ograniczeń, do których
raport wraca w sekcji~\ref{sec:limitations}.

Skoring działa w jednym przejściu modelu przez tekst złożony
z promptu i wygenerowanej kontynuacji. Tokeny promptu są maskowane,
więc do metryk wliczają się tylko te tokeny, które model sam
wygenerował - mierzona jest jego pewność co do własnego rozumowania,
a nie znajomość samego polecenia. Implementacja znajduje się w pliku
\texttt{experiments/perplexity.py} (funkcja \texttt{score\_continuation}).

\subsection{Dwa warianty: CoT i agent ReAct}
\label{sec:settings}

Pierwszy wariant to czysty chain-of-thought: model dostaje prompt
zachęcający do myślenia krok po kroku i generuje całą odpowiedź
w jednym przebiegu, bez narzędzi.

Drugi wariant to pętla ReAct~\cite{yao2023react}, w której model
na przemian produkuje segmenty oznaczone \texttt{Thought:},
\texttt{Action:}, \texttt{Observation:} i \texttt{Answer:}.
Jedynym dostępnym narzędziem jest kalkulator, a liczba wywołań
w jednym pytaniu ograniczona jest do pięciu.

Te same 60 pytań i ten sam model (Qwen2.5-1.5B) w obu wariantach
pozwalają czysto porównać, czy włożenie modelu w pętlę agentową
zmienia samą relację między PPL a poprawnością - nawet jeśli
dokładność spada lub rośnie.

\subsection{Metryki}
\label{sec:metryki}

Rozkład prawdopodobieństw tokenów sprowadzany jest do czterech
skalarów. Pierwszy z nich, \textit{perplexity}, dane jest
równaniem~\eqref{eq:ppl} i mówi, jak bardzo średnio model był
zaskoczony własnymi tokenami. \textit{Średnia entropia tokenowa}
$-\sum_v p(v)\log p(v)$, uśredniona po wszystkich pozycjach
sekwencji, korzysta z całego rozkładu, a nie tylko z tokena, który
został wybrany - dzięki temu reaguje także wtedy, gdy top-1 token
jest pewny, ale za nim stoi ,,gęsta'' alternatywa.
\textit{Top-1 margin} to różnica $p(\text{top-1}) - p(\text{top-2})$,
czyli to, jak bardzo wybrany token wygrał z najbliższym konkurentem.
\textit{Top-1 mass} to po prostu $p(\text{top-1})$.

Do oceny kalibracji wykorzystano \emph{Expected Calibration Error}
(ECE), liczone na binach o równej liczbie obserwacji
($n_{\text{bins}}=10$). Pewność modelu definiowana jest tu jako
$1/\mathrm{PPL}$. Niższe ECE oznacza lepszą zgodność tego, jak
pewny był model, z tym, jak często rzeczywiście miał rację.
